% 同一教学场景中的图像级、关键点与边界标注；斜线区表示未检查区域。
\begin{tikzpicture}[x={\dimexpr\linewidth/182\relax},y=1mm,scale=.98,transform shape,font=\figurefont\footnotesize,text=BookInk]
  \tikzset{
    scene/.style={draw=BookRule,fill=BookPaper,line width=.8pt},
    object/.style={line width=1pt},
    selected boundary/.style={draw=FigureModel,line width=1.3pt},
    panel title/.style={font=\figurefont\bfseries\footnotesize,anchor=south},
    output/.style={draw=BookRule,fill=white,minimum width=48mm,minimum height=8mm,align=center}
  }
  \foreach \shift/\panel/\title in {0/a/图像级类别,61/b/对象关键点,122/c/对象边界}{
    \begin{scope}[xshift=\shift mm]
      \node[panel title] at (26,49) {(\panel)\quad\title};
      \path[scene] (1,8) rectangle (51,43);
      \draw[object,draw=BookBlue,fill=FigureInput!18] (14,27) circle (6);
      \draw[object,draw=BookAmber,fill=FigureLoss!18] (25,13)--(37,13)--(31,27)--cycle;
      \draw[object,draw=BookTeal,fill=FigureModel!18] (39,30) rectangle (48,38);
    \end{scope}
  }
  \begin{scope}
    \node[output] at (26,3) {输出：场景含目标对象};
  \end{scope}
  \begin{scope}[xshift=61mm]
    \begin{scope}
      \clip (25,8) rectangle (51,43);
      \foreach \d in {-20,-15,...,70}{\draw[BookRule,line width=.55pt] (\d,8)--({\d+35},43);}
    \end{scope}
    \draw[BookMuted,dashed,line width=.8pt] (25,8) rectangle (51,43);
    \fill[FigureModel] (14,27) circle (1.5);
    \draw[white,line width=.7pt] (12.6,27)--(15.4,27) (14,25.6)--(14,28.4);
    \node[output] at (26,3) {输出：对象甲的关键点};
    \node[anchor=north east,text=BookMuted] at (49.5,41.5) {未检查};
  \end{scope}
  \begin{scope}[xshift=122mm]
    \begin{scope}
      \clip (38,8) rectangle (51,43);
      \foreach \d in {-20,-15,...,70}{\draw[BookRule,line width=.55pt] (\d,8)--({\d+35},43);}
    \end{scope}
    \draw[BookMuted,dashed,line width=.8pt] (38,8) rectangle (51,43);
    \draw[selected boundary] (7,20) rectangle (21,34);
    \draw[selected boundary] (23,11) rectangle (38,29);
    \node[output] at (26,3) {输出：已检查对象甲、乙的边界};
    \node[anchor=north east,text=BookMuted] at (49.5,41.5) {未检查};
  \end{scope}
\end{tikzpicture}
